\section{Метод конечных элементов. Лагранжевы элементы}
\label{sec:fem-lagrange}

\subsection{Идея метода конечных элементов}

Метод конечных элементов (МКЭ) применяется к краевым задачам для
дифференциальных уравнений. Область разбивается на простые подобласти ---
конечные элементы, а искомая функция приближается кусочно-полиномиальной
функцией. Коэффициенты определяются из слабой, или вариационной, постановки.

Рассмотрим модельную задачу Пуассона:
\[
    -\nabla\cdot(k\nabla u)=f\quad\text{в }\Omega,
    \qquad
    u=0\quad\text{на }\Gamma_D,
\]
с возможным условием Неймана
\[
    k\frac{\partial u}{\partial n}=g\quad\text{на }\Gamma_N.
\]

Умножим уравнение на тестовую функцию $v$, проинтегрируем по области и применим
формулу Грина:
\[
    \int_{\Omega} k\nabla u\cdot\nabla v\,dx
    =
    \int_{\Omega}fv\,dx+
    \int_{\Gamma_N}gv\,ds.
\]
Слабая постановка: найти $u\in V$, где
\[
    V=\{v\in H^1(\Omega):v|_{\Gamma_D}=0\},
\]
такое, что
\[
    a(u,v)=\ell(v)
    \qquad\forall v\in V.
\]

\subsection{Корректность слабой постановки: теорема Лакса--Мильграма}

Пусть $V$ --- гильбертово пространство, билинейная форма $a(\cdot,\cdot)$
непрерывна и коэрцитивна:
\[
    |a(u,v)|\le M\|u\|_V\|v\|_V,
    \qquad
    a(v,v)\ge \alpha\|v\|_V^2,
    \qquad \alpha>0,
\]
а $\ell$ --- непрерывный линейный функционал на $V$. Тогда по
\textbf{теореме Лакса--Мильграма} существует единственное $u\in V$, для которого
\[
    a(u,v)=\ell(v)\qquad\forall v\in V,
\]
и выполняется оценка $\|u\|_V\le \|\ell\|_{V^*}/\alpha$.
Для модельной эллиптической задачи коэрцитивность обычно следует из
положительности коэффициента $k$ и неравенства Пуанкаре.

\subsection{Галёркинская аппроксимация}

Выбирается конечномерное пространство $V_h\subset V$ с базисом
$\{\varphi_1,\ldots,\varphi_N\}$. Приближённое решение имеет вид
\[
    u_h=\sum_{j=1}^{N}U_j\varphi_j.
\]
Требование
\[
    a(u_h,v_h)=\ell(v_h)
    \qquad\forall v_h\in V_h
\]
приводит к системе
\[
    KU=F,
\]
где
\[
    K_{ij}=a(\varphi_j,\varphi_i),
    \qquad
    F_i=\ell(\varphi_i).
\]
Матрица $K$ называется матрицей жёсткости.

\subsection{Лагранжевы конечные элементы}

Лагранжев элемент задаётся тройкой $(K,P,\Sigma)$:
\begin{itemize}
    \item $K$ --- геометрический элемент;
    \item $P$ --- пространство полиномов;
    \item $\Sigma$ --- набор степеней свободы, обычно значений функции в узлах.
\end{itemize}

Базисные функции Лагранжа удовлетворяют условию
\[
    \varphi_i(x_j)=\delta_{ij}.
\]
Поэтому коэффициент $U_i$ равен узловому значению приближённого решения.

Для одномерного линейного элемента $K=[x_1,x_2]$:
\[
    \varphi_1(x)=\frac{x_2-x}{h},
    \qquad
    \varphi_2(x)=\frac{x-x_1}{h},
    \qquad h=x_2-x_1.
\]
На элементе
\[
    u_h(x)=U_1\varphi_1(x)+U_2\varphi_2(x).
\]

\subsection{Линейный треугольный элемент}

На треугольнике используются полиномы первой степени
\[
    P_1(K)=\operatorname{span}\{1,x,y\}.
\]
Степени свободы --- значения в трёх вершинах. Через барицентрические координаты
$\lambda_1,\lambda_2,\lambda_3$ базис записывается как
\[
    \varphi_i=\lambda_i,
    \qquad
    \lambda_1+\lambda_2+\lambda_3=1.
\]
Градиенты базисных функций постоянны внутри элемента, поэтому градиент линейного
конечно-элементного решения кусочно постоянен.

Квадратичный треугольный элемент имеет шесть узлов: три вершины и три середины
рёбер. Его базис:
\[
\begin{aligned}
\varphi_i&=\lambda_i(2\lambda_i-1),\qquad i=1,2,3,\\
\varphi_{ij}&=4\lambda_i\lambda_j.
\end{aligned}
\]

\subsection{Локальные матрицы и сборка}

На каждом элементе $K_e$ вычисляются локальные матрица и вектор:
\[
    K_{ij}^{(e)}=
    \int_{K_e}k\nabla\varphi_j^{(e)}\cdot\nabla\varphi_i^{(e)}\,dx,
\]
\[
    F_i^{(e)}=
    \int_{K_e}f\varphi_i^{(e)}\,dx
    +\int_{\partial K_e\cap\Gamma_N}g\varphi_i^{(e)}\,ds.
\]
Затем локальные вклады суммируются в общую разреженную систему по глобальным
номерам степеней свободы. Этот процесс называется сборкой.

Для одномерного элемента при постоянном $k$ локальная матрица жёсткости равна
\[
    K^{(e)}=\frac{k}{h}
    \begin{pmatrix}
    1&-1\\
    -1&1
    \end{pmatrix}.
\]

\subsection{Отображение с эталонного элемента}

Вычисления удобно выполнять на эталонном элементе $\widehat K$. Геометрическое
отображение
\[
    x=F_e(\widehat x)
\]
переносит базисные функции и интегралы на физический элемент. Для
изопараметрических элементов геометрия и решение аппроксимируются одними и теми
же функциями формы.

При замене переменных учитываются якобиан и преобразование градиента:
\[
    \nabla_x\varphi=J^{-T}\nabla_{\widehat x}\widehat\varphi,
    \qquad
    dx=|\det J|\,d\widehat x.
\]

\subsection{Граничные условия}

Условия Неймана входят в правую часть слабой постановки естественным образом.
Условия Дирихле являются существенными: соответствующие степени свободы задаются
явно, после чего система модифицируется исключением известных значений или
другим эквивалентным способом.

\subsection{Сходимость и оценка погрешности}

При непрерывности $a$ с константой $M$ и коэрцитивности с константой $\alpha$
\textbf{лемма Сеа} даёт
\[
    \|u-u_h\|_V
    \leq \frac{M}{\alpha}
    \inf_{v_h\in V_h}\|u-v_h\|_V.
\]
То есть ошибка МКЭ в энергетической норме сравнима с наилучшей ошибкой
аппроксимации выбранным конечномерным пространством.

Для Лагранжевых элементов степени $p$ на регулярном семействе сеток при
достаточной гладкости решения типична энергетическая оценка
\[
    \|u-u_h\|_{H^1(\Omega)}\leq C h^p|u|_{H^{p+1}(\Omega)},
\]
\[
    \|u-u_h\|_{L^2(\Omega)}\leq C h^{p+1}|u|_{H^{p+1}(\Omega)}.
\]
Последняя, более сильная $L^2$-оценка требует дополнительной регулярности
сопряжённой эллиптической задачи (аргумент Обена--Нитше).

\subsection{Принцип Ритца и энергетическая интерпретация}

Если билинейная форма $a(\cdot,\cdot)$ симметрична, слабая задача
\[
    a(u,v)=\ell(v)\qquad \forall v\in V
\]
эквивалентна минимизации функционала энергии
\[
    J(v)=\frac12 a(v,v)-\ell(v).
\]
Действительно, первая вариация равна
\[
    J'(u)[w]=a(u,w)-\ell(w),
\]
и стационарность при всех $w\in V$ даёт слабое уравнение. При коэрцитивности форма
задаёт строго выпуклый функционал, поэтому стационарная точка является единственным
минимумом.

Метод Ритца ограничивает минимизацию подпространством $V_h\subset V$, а метод Галёркина
требует ортогональности невязки ко всем тестовым функциям из $V_h$. Для симметричной
задачи эти два подхода совпадают.

\subsection{Лемма Сеа: схема доказательства}

Пусть
\[
    |a(w,v)|\le M\|w\|_V\|v\|_V,
    \qquad
    a(v,v)\ge \alpha\|v\|_V^2.
\]
Из галёркинской ортогональности
\[
    a(u-u_h,v_h)=0\qquad\forall v_h\in V_h
\]
для произвольного $w_h\in V_h$ имеем
\[
\begin{aligned}
    \alpha\|u-u_h\|_V^2
    &\le a(u-u_h,u-u_h)\\
    &=a(u-u_h,u-w_h)\\
    &\le M\|u-u_h\|_V\|u-w_h\|_V.
\end{aligned}
\]
Отсюда
\[
    \|u-u_h\|_V
    \le \frac{M}{\alpha}
    \inf_{w_h\in V_h}\|u-w_h\|_V.
\]
То есть МКЭ квазиоптимален: ошибка определяется качеством аппроксимации выбранного
конечно-элементного пространства.

\subsection{Интерполяционные оценки для лагранжевых элементов}

Пусть $I_hu$ --- лагранжев интерполянт степени $p$ на регулярном семействе сеток.
При достаточной гладкости $u$ стандартная локальная оценка имеет вид
\[
    |u-I_hu|_{H^m(K)}
    \le C h_K^{p+1-m}|u|_{H^{p+1}(K)},
    \qquad 0\le m\le p+1.
\]
После суммирования по элементам и применения леммы Сеа получаем типичную энергетическую
оценку
\[
    \|u-u_h\|_{H^1(\Omega)}
    \le Ch^p|u|_{H^{p+1}(\Omega)}.
\]
Улучшенная $L^2$-оценка порядка $h^{p+1}$ требует дополнительной регулярности
сопряжённой эллиптической задачи (аргумент Обена--Ницше).

\subsection{Структура конечно-элементной матрицы}

Для эллиптической задачи локальность функций формы приводит к разреженной матрице
жёсткости: элемент $K_{ij}$ ненулевой только тогда, когда носители $\phi_i$ и $\phi_j$
пересекаются. При симметричной коэрцитивной постановке матрица после наложения
существенных граничных условий симметрична и положительно определена.

При измельчении сетки число неизвестных растёт, а обусловленность матрицы обычно
ухудшается. Поэтому для больших систем применяют итерационные методы и
предобуславливатели; геометрические и алгебраические многосеточные методы особенно
естественны для конечно-элементных задач.

\subsection{Что важно сказать на экзамене}

Необходимо показать цепочку: сильная постановка --- интегрирование по частям ---
слабая постановка --- теорема Лакса--Мильграма --- конечномерное пространство ---
система $KU=F$ --- лемма Сеа.
Лагранжевы элементы используют значения функции в узлах и обеспечивают
$C^0$-непрерывность между элементами.

\newpage